\documentclass[preprint]{aastex631}

\begin{document}

\title{Validation of SWRF-Impute: MCAR, MAR, MNAR, and Realistic Stellar Host Testing}
\author{Agastya Gaur}
\noaffiliation

\section{Introduction}

Validation of an imputation algorithm is essential to assess its ability to reconstruct missing values accurately while preserving statistical properties of the dataset. The reliability of imputation depends critically on the missing data mechanism. This document presents validation methodology and completed results for the SWRF-Impute imputation algorithm under four distinct missing data mechanisms: Missing Completely At Random (MCAR), Missing At Random (MAR), Missing Not At Random (MNAR), and Realistic Stellar Host Missingness (RSHM).

The first three mechanisms (MCAR, MAR, MNAR) are standard benchmarks that characterize the full range of missingness structures in the statistical literature. RSHM is a domain-specific mechanism introduced in this work to simulate the detection-driven selection effects of exoplanet surveys: stars hosting planets that are unlikely to be detected by radial velocity or transit methods are systematically underrepresented in confirmed-host catalogs, producing MNAR-like missingness whose structure is physically motivated rather than arbitrarily constructed.

All validation studies are conducted on a synthetic stellar host population generated to match the distributional properties of real exoplanet host stars while providing fully observed ground truth. This design enables rigorous evaluation of covariance structure recovery under each mechanism, including RSHM, where real data lacks a complete-case reference.

\section{Data and Features}

The validation studies employ a synthetic stellar host population. The synthetic population is constructed to reflect the distributional properties of real exoplanet host stars while providing a fully observed ground truth matrix with no pre-existing missingness. This allows rigorous evaluation across all four mechanisms (MCAR, MAR, MNAR, RSHM) without the confound of pre-existing real-world missingness in the baseline.

The data matrix is $\mathbf{X} \in \mathbb{R}^{N \times P}$ with $N = 500$ synthetic stars and $P = 9$ features. The features are listed in \autoref{tab:data_parameters}.

\begin{table}
  \centering
  \begin{tabular}{lll}
    \hline
    \textbf{Parameter} & \textbf{Description} & \textbf{Source} \\
    \hline
    \verb|sy_pnum| & Number of Planets & Poisson draw, metallicity-coupled \\
    \verb|sy_snum| & Number of Stars in System & Poisson draw \\
    \verb|st_teff| & Stellar Effective Temperature [K] & Gaia DR3 empirical sample \\
    \verb|st_rad|  & Stellar Radius [R$_\odot$] & Gaia DR3 empirical sample \\
    \verb|st_mass| & Stellar Mass [M$_\odot$] & Gaia DR3 empirical sample \\
    \verb|st_met_FeH| & Stellar Metallicity [Fe/H] [dex] & Gaia DR3, planet-count scaled \\
    \verb|st_lum|  & Stellar Luminosity [W] & Derived: $4\pi R^2 \sigma T_{\text{eff}}^4$ \\
    \verb|st_logg| & Stellar Surface Gravity [log$_{10}$(cm/s$^2$)] & Derived: $\log_{10}(g_\odot M/R^2)$ \\
    \verb|st_age|  & Stellar Age [Gyr] & Gaia DR3 empirical sample \\
    \hline
  \end{tabular}
  \caption{Features of the synthetic stellar host population used in all validation studies. Stellar parameters are drawn from Gaia DR3 empirical distributions. Luminosity and surface gravity are computed analytically from sampled radius, mass, and temperature. A planet-metallicity coupling is imposed: \texttt{st\_met\_FeH} is increased by 20\% of its absolute value per additional planet, consistent with the known giant-planet--metallicity correlation.}
  \label{tab:data_parameters}
\end{table}

\subsection{Synthetic Population Design}

The synthetic population is fully observed by construction, providing clean ground truth for all validation experiments. Stellar parameters (temperature, radius, mass, metallicity, age) are drawn independently from Gaia DR3 empirical distributions. Planet count \texttt{sy\_pnum} follows a shifted Poisson distribution ($\lambda = 1.5$, minimum 1), reduced by the number of companion stars to maintain physical plausibility. A planet-metallicity coupling is imposed:

\begin{equation}
  \texttt{st\_met\_FeH} \leftarrow \texttt{st\_met\_FeH} + 0.2 \times |\texttt{st\_met\_FeH}| \times \texttt{sy\_pnum}
\end{equation}

This introduces a structural correlation between \verb|sy_pnum| and \verb|st_met_FeH|, consistent with the well-established giant-planet--metallicity correlation \citep{johnsonGiantPlanetOccurrence2010,gonzalezStellarMetallicityGiant1997}. Luminosity and surface gravity are computed analytically and therefore carry the correlations of their parent parameters (\texttt{st\_rad}, \texttt{st\_mass}, \texttt{st\_teff}) into the covariance structure.

The resulting $15{,}000 \times 9$ matrix provides a rich, physically motivated covariance structure against which imputation methods can be benchmarked.

\subsection{Physical Bounds}

SWRF-Impute is initialized via a bounds list: a sequence of $P$ pairs $(l_p,\, u_p)$ that specify the physically meaningful lower and upper limit for each feature. Either bound may be absent.

Missing values in feature $p$ are initialized from $\mathrm{Uniform}(l_p, u_p)$ when both bounds are present, and from the observed column mean (clamped to any present bound) otherwise. During Gibbs sampling, the truncated-normal draw for each feature is clipped to $(l_p, u_p)$; absent bounds leave the corresponding tail unconstrained.

For the validation studies, bounds will be set to physically motivated ranges for each stellar parameter (e.g.\ positive-only constraints on mass and radius, temperature limits consistent with main-sequence stars). Features for which no meaningful constraint exists are left unbounded.

\section{Evaluation Metrics}

Two metrics assess imputation quality:

\subsection{Imputation Accuracy: RMSE on Masked Entries}

The primary metric quantifies the pointwise accuracy of imputed values. Given the original complete data $\mathbf{X}_{orig}$, imputed data $\mathbf{X}_{imp}$, and binary mask $\mathbf{M}$ indicating missing locations, we compute:

\begin{equation}
  \text{RMSE}_{\text{masked}} = \sqrt{\frac{1}{\sum_{i,j} \mathbf{M}_{i,j}} \sum_{i,j} \mathbf{M}_{i,j} (\mathbf{X}_{orig,i,j} - \mathbf{X}_{imp,i,j})^2}
\end{equation}

This metric measures how close imputed values are to the ground truth on missing entries only, capturing the algorithm's pointwise predictive accuracy. Lower values indicate better imputation quality.

\subsection{Pattern Preservation: Eigenvalue Spectrum Distance}

The second metric assesses whether imputation preserves the statistical structure and correlations in the data. After mean-centering and standardizing both the original and imputed data, Principal Component Analysis (PCA) is performed to obtain eigenvalue spectra $\mathbf{\lambda}_{true}$ and $\mathbf{\lambda}_{imp}$ representing the variance structure.

We quantify the distance between spectra using logarithmic $L^2$ distance:

\begin{equation}
  d_{\text{eig}} = \sqrt{\sum_{k=1}^{P} (\log(\lambda_{true,k} + \epsilon) - \log(\lambda_{imp,k} + \epsilon))^2}
\end{equation}

where $\epsilon = 10^{-8}$ provides numerical stability. This log-scale distance is appropriate for eigenvalues spanning multiple orders of magnitude and emphasizes relative differences. Lower values indicate that the imputation preserves the correlational structure of the data, which is critical for downstream analysis.

\subsection{Uncertainty Calibration: IQR Coverage}

RMSE and eigenvalue distance both evaluate the \emph{point estimate} produced by pooling an imputation ensemble. Neither captures whether the ensemble itself reflects well-calibrated uncertainty about the missing values. To probe this, we report a complementary metric that is only defined for stochastic methods.

Let $\mathcal{E} = \{\mathbf{X}_{\text{imp}}^{(k)}\}_{k=1}^{K}$ be the ensemble of $K$ imputed matrices produced by a stochastic method. For each cell $(i,j)$ we compute the 25th and 75th percentiles across the ensemble:

\begin{equation}
  Q_{25}(i,j) = \text{Percentile}_{25}\bigl(\{\mathbf{X}_{\text{imp}}^{(k)}(i,j)\}_{k=1}^{K}\bigr), \quad
  Q_{75}(i,j) = \text{Percentile}_{75}\bigl(\{\mathbf{X}_{\text{imp}}^{(k)}(i,j)\}_{k=1}^{K}\bigr)
\end{equation}

The IQR coverage on masked entries is then the fraction of ground-truth values that fall inside the cell-wise interquartile range:

\begin{equation}
  C_{\text{IQR}} = \frac{1}{\sum_{i,j} \mathbf{M}_{i,j}} \sum_{i,j} \mathbf{M}_{i,j} \cdot \mathbf{1}\bigl[Q_{25}(i,j) \leq \mathbf{X}_{orig,i,j} \leq Q_{75}(i,j)\bigr]
\end{equation}

A well-calibrated imputer draws values that spread around the truth in proportion to its uncertainty; because an IQR spans the middle 50\% of a distribution by construction, such an imputer should have $C_{\text{IQR}} \approx 0.5$. Values below $0.5$ indicate under-dispersed (over-confident) draws --- the ensemble concentrates too tightly to cover the truth half the time. Values above $0.5$ indicate over-dispersed (under-confident) draws. This metric directly probes one of SWRF-Impute's central claims: that its multi-draw Gibbs-style output produces calibrated, rather than merely plausible, uncertainty estimates. The IQR (as opposed to a wider interval such as 90\%) is chosen because with $K = 35$ draws per ensemble the outer quantiles are noisy estimates of the tails.

The metric is only defined for stochastic methods. For MICE and MissForest the ensemble is the set of $n_{\text{runs}}$ independently sampled imputations; for SWRF-Impute it is the post-burn-in Gibbs chain.

\section{Validation Methodology}

Across all four missing data mechanisms (MCAR, MAR, MNAR, and RSHM), the validation pipeline is identical except for mask generation.

\subsection{Initialization}

All tests begin with the same setup. Starting from the fully observed synthetic matrix $\mathbf{X} \in \mathbb{R}^{N \times P}$ with $N = 15{,}000$ and $P = 9$, we systematically introduce missingness:

\begin{equation}
  n_{\text{missing}} = \lfloor \text{prop\_missing} \times (N \times P) \rfloor
\end{equation}

where $\text{prop\_missing}$ ranges from $0.1$ (10\%) to $0.6$ (60\%) in increments of $0.1$. The missingness is applied via mechanism-specific mask generation (see \S\ref{sec:mask-generation} for details on MCAR, MAR, MNAR, and RSHM mask creation). Because the synthetic matrix is fully observed, ground truth is available for every masked entry under every mechanism, including RSHM.

\subsection{Imputation Algorithm Comparison}

The following six algorithms are evaluated identically for all three mechanisms:

\begin{enumerate}
  \item \textbf{Mean Imputation} - Replaces missing values with the mean of observed values in each feature.
  \item \textbf{Median Imputation} - Replaces missing values with the median of observed values in each feature.
  \item \textbf{$k$-Nearest Neighbors (KNN)} - Imputes using the average of the 5 nearest complete neighbors in standardized feature space.
  \item \textbf{MICE} - Multivariate Imputation by Chained Equations with maximum 10 iterations and posterior sampling to introduce stochasticity.
  \item \textbf{MissForest} - Iterative imputation using Random Forest regressors (30 trees, 10 iterations), unweighted.
  \item \textbf{SWRF-Impute} - The proposed method: weighted stochastic imputation via random forests (30 trees, 10 iterations) with truncated normal sampling.
\end{enumerate}

\subsection{Stochastic Method Handling}

All stochastic methods (MICE, MissForest, and SWRF-Impute) are evaluated on a pooled point estimate: multiple imputed matrices are generated and then averaged cell-wise into a single matrix before computing RMSE and eigenvalue distance. Evaluating all stochastic methods on the pooled matrix, rather than averaging per-run metrics, places them on equal footing and is consistent with the Rubin's rules approach to multiple imputation.

\begin{enumerate}
  \item \textbf{MICE}: executed $n_{\text{runs}} = 35$ independent times per missingness level, each with a distinct seed drawn from a seeded random generator. The 35 resulting imputed matrices are averaged cell-wise to form the pooled point estimate used for evaluation.
  \item \textbf{MissForest}: executed $n_{\text{runs}} = 35$ independent times per missingness level, each with a distinct seed, and the 35 imputed matrices are averaged cell-wise. Unlike simple deterministic use of MissForest (single run with a fixed seed), this accounts for the variance introduced by Random Forest bootstrap sampling.
  \item \textbf{SWRF-Impute}: a single Gibbs-style chain is run with $n_{\text{iters}} = 15$ iterations and $\text{draws\_per\_iter} = 3$, producing 45 total draws. After discarding a burn-in of $b = 10$ draws, the remaining 35 post-burn-in draws are averaged cell-wise to form the pooled point estimate.
\end{enumerate}

\subsection{Reproducibility}

All mask generation and stochastic imputation is driven by a single seeded \texttt{numpy.random.Generator} (via \texttt{numpy.random.default\_rng}). A master seed initialises the generator; independent child generators for each (iteration, missingness level) pair are derived by drawing 32-bit integer seeds from the master. This ensures all experiments are fully reproducible while keeping individual mask realisations statistically independent of one another.

The seed thread runs through the entire pipeline. SWRF-Impute's internal random operations---initialization draws (uniform within bounds, or mean fill), Random Forest bootstrap sampling, and truncated-normal draws at each Gibbs step---all accept the child generator directly. MICE and MissForest receive integer seeds derived from the same master generator for each of their $n_{\text{runs}}$ calls.

\subsection{Result Aggregation}

A single full test consists of:
\begin{itemize}
  \item Evaluation at 6 missingness levels: $\{0.1, 0.2, \ldots, 0.6\}$
  \item 6 imputation algorithms
\end{itemize}

The test can be repeated across multiple iterations ($n_{\text{tot}}$ trials) to assess stability. Final performance metrics are averaged across all iterations, producing two result dictionaries indexed by mechanism and missingness level.

\begin{equation}
  \text{RMSE}_{\text{results}} = \{\text{algorithm} : [\text{RMSE}_{0.1}, \ldots, \text{RMSE}_{0.9}]\}
\end{equation}

\begin{equation}
  \text{Eigval}_{\text{results}} = \{\text{algorithm} : [d_{\text{eig},0.1}, \ldots, d_{\text{eig},0.9}]\}
\end{equation}

\section{Mask Generation Methods}
\label{sec:mask-generation}

The primary difference across MCAR, MAR, and MNAR testing is the mask generation strategy. All other pipeline components (initialization, imputation, aggregation) are identical.

\subsection{MCAR}

In MCAR data, missing data is completely random. To simulate this for each missingness level, a single random mask is generated for each full test. The mask is created by:

\begin{enumerate}
  \item Flattening the complete data matrix to a vector of length $N \times P$
  \item Randomly sampling $n_{\text{missing}}$ indices without replacement
  \item Converting indices back to (row, column) pairs via $\text{unravel\_index}()$
  \item Setting masked entries to NaN to create $\mathbf{X}_{\text{masked}}^{\text{MCAR}}$
\end{enumerate}

\subsection{MAR}

In MAR data, the probability that a value is missing depends on observed variables. We generate MAR missingness by selecting a random feature as a trigger variable and sampling cells to mask from all non-trigger columns with row probability proportional to the extremity of the trigger feature. Unlike a hard row cutoff, this produces a smooth gradient: the most extreme observations are the most likely to contribute missing values, but no observation is categorically excluded. This method is loosely based on the MAR generation method used by \citet{toyeBenchmarkingMissingData2025}.

\begin{enumerate}
  \item \textbf{Select trigger feature}: Randomly choose one feature index $f_{\text{trigger}} \in \{1, \ldots, P\}$.
  \item \textbf{Compute row weights}: Standardize the trigger column to z-scores $z_i = (x_{i,f_{\text{trigger}}} - \mu) / \sigma$ and set row weight $w_i = |z_i|$. Normalize so that $\sum_i w_i = 1$.
  \item \textbf{Build eligible cell set}: Let $\mathcal{C} = \{(i, j) : j \neq f_{\text{trigger}},\; i = 1, \ldots, N\}$ be all non-trigger cells. Assign each cell $(i, j) \in \mathcal{C}$ the weight $w_i$ (row weight, shared across all non-trigger columns in that row), then renormalize over all cells.
  \item \textbf{Sample without replacement}: Draw exactly $n_{\text{missing}} = \lfloor \text{prop\_missing} \times N \times P \rfloor$ cells from $\mathcal{C}$ without replacement according to the cell weights and set them to NaN.
\end{enumerate}

The trigger column remains fully observed. All other columns participate in missingness with per-row probability proportional to $|z_i|$, so the mechanism is ignorable (MAR) given the trigger: an imputer with access to the trigger can in principle model the selection process.

\subsection{MNAR}

In MNAR data, the probability that a value is missing depends on the unobserved value itself. We employ a \emph{stochastic tail masking} strategy: each cell is assigned a missing probability based on its z-score magnitude in its feature, and cells are then sampled stochastically to reach the target proportion. This is physically motivated for astrophysical data---stars with extreme stellar parameters (very high or very low temperature, mass, metallicity, etc.) are precisely the objects for which accurate measurements are hardest to obtain, and therefore most likely to be absent from catalogs.

\begin{enumerate}
  \item \textbf{Compute per-cell z-scores}: For each feature $j \in \{1, \ldots, P\}$, compute $z_{ij} = (x_{ij} - \mu_j) / \sigma_j$ for all $i$, and take the absolute value: $|z_{ij}|$.
  \item \textbf{Assign missing probabilities}: For each cell $(i,j)$, assign missing probability $p_{ij} = |z_{ij}|^2 / \sum_{i,j} |z_{ij}|^2$ (normalized over all cells). The quadratic dependence emphasizes extreme values.
  \item \textbf{Stochastically sample missing cells}: Draw $n_{\text{miss}} = \max\bigl(1,\; \text{round}(\text{prop\_missing} \times N \times P)\bigr)$ cells without replacement according to their probabilities $p_{ij}$. Set these cells to NaN.
\end{enumerate}

The total number of masked cells is $n_{\text{miss}} \approx \text{prop\_missing} \times N \times P$, matching the global target proportion. Because missingness across all cells is driven by the unobserved values themselves (via the z-score weighting), the mechanism is non-ignorable (MNAR): no subset of the remaining observed data can fully explain which entries are absent. The stochastic sampling ensures that mask generation varies between runs, providing a more realistic and stable evaluation across different random seeds.

\subsection{Realistic Stellar Host Missingness (RSHM)}

\textit{[Placeholder — implementation in progress.]}

RSHM simulates the detection-driven selection effects present in real exoplanet host star catalogs. Stars hosting planets that fall below survey detection thresholds (due to long orbital periods, low planet-to-star radius ratios, or unfavorable stellar noise properties) are systematically underrepresented in confirmed-host samples. This produces MNAR-like missingness whose structure is physically motivated by the coupling between stellar parameters (\texttt{st\_teff}, \texttt{st\_mass}, \texttt{st\_met\_FeH}) and planet detectability.

The mechanism will be parameterized by a detection probability function $p_{\text{detect}}(i)$ that depends on each star's stellar properties and planet count, with stars assigned to missing status with probability $1 - p_{\text{detect}}(i)$ across multiple features.

\bibliography{references}

\end{document}
